\documentclass[../../main]{subfiles}

\begin{document}

\section{Operações entre Conjuntos}
\label{sec:matematica:teoria-conjuntos:operacoes-entre-conjuntos}

\begin{defn}[União]
  \label{def:matematica:teoria-conjuntos:operacoes:uniao}

  Sejam (A) e (B) conjuntos.

  A união de (A) e (B), denotada por

  [
    A \cup B,
  ]

  é o conjunto definido por

  [
    A \cup B
    =
    {
      x
      \mid
      x \in A
      \lor
      x \in B
    }.
  ]

\end{defn}

\begin{defn}[Interseção]
  \label{def:matematica:teoria-conjuntos:operacoes:intersecao}

  Sejam (A) e (B) conjuntos.

  A interseção de (A) e (B), denotada por

  [
    A \cap B,
  ]

  é o conjunto definido por

  [
    A \cap B
    =
    {
      x
      \mid
      x \in A
      \land
      x \in B
    }.
  ]

\end{defn}

\begin{defn}[Diferença]
  \label{def:matematica:teoria-conjuntos:operacoes:diferenca}

  Sejam (A) e (B) conjuntos.

  A diferença entre (A) e (B), denotada por

  [
    A \setminus B,
  ]

  é definida por

  [
    A \setminus B
    =
    {
      x
      \mid
      x \in A
      \land
      x \notin B
    }.
  ]

\end{defn}

\begin{defn}[Diferença Simétrica]
  \label{def:matematica:teoria-conjuntos:operacoes:diferenca-simetrica}

  Sejam (A) e (B) conjuntos.

  A diferença simétrica de (A) e (B) é definida por

  [
    A \triangle B
    =
    (A \setminus B)
    \cup
    (B \setminus A).
  ]

\end{defn}

\begin{defn}[Complemento]
  \label{def:matematica:teoria-conjuntos:operacoes:complemento}

  Seja (U) um conjunto universal e seja (A \subseteq U).

  O complemento de (A) em (U), denotado por

  [
    A^{c},
  ]

  é definido por

  [
    A^{c}
    =
    U \setminus A.
  ]

\end{defn}

\subsection{Propriedades Fundamentais}

\begin{prop}[Idempotência da União]

  Para todo conjunto (A),

  [
    A \cup A = A.
  ]

\end{prop}

\begin{proof}

  Para todo elemento (x),

  [
    x \in A \cup A
    \Leftrightarrow
    (x \in A \lor x \in A).
  ]

  Pela idempotência da disjunção,

  [
    x \in A \cup A
    \Leftrightarrow
    x \in A.
  ]

  Pelo princípio da extensionalidade,

  [
    A \cup A = A.
  ]

\end{proof}

\begin{prop}[Idempotência da Interseção]

  Para todo conjunto (A),

  [
    A \cap A = A.
  ]

\end{prop}

\begin{proof}

  Análoga à demonstração anterior.

\end{proof}

\begin{prop}[Comutatividade da União]

  Para quaisquer conjuntos (A) e (B),

  [
    A \cup B = B \cup A.
  ]

\end{prop}

\begin{proof}

  Seja (x).

  Então

  [
    x \in A \cup B
    \Leftrightarrow
    (x \in A \lor x \in B).
  ]

  Pela comutatividade da disjunção,

  [
    x \in A \cup B
    \Leftrightarrow
    (x \in B \lor x \in A).
  ]

  Portanto

  [
    x \in A \cup B
    \Leftrightarrow
    x \in B \cup A.
  ]

  Pela extensionalidade,

  [
    A \cup B = B \cup A.
  ]

\end{proof}

\begin{prop}[Comutatividade da Interseção]

  Para quaisquer conjuntos (A) e (B),

  [
    A \cap B = B \cap A.
  ]

\end{prop}

\begin{proof}

  Análoga à demonstração anterior.

\end{proof}

\begin{prop}[Associatividade da União]

  Para quaisquer conjuntos (A), (B) e (C),

  [
    (A \cup B)\cup C
    =
    A\cup(B\cup C).
  ]

\end{prop}

\begin{prop}[Associatividade da Interseção]

  Para quaisquer conjuntos (A), (B) e (C),

  [
    (A \cap B)\cap C
    =
    A\cap(B\cap C).
  ]

\end{prop}

\begin{prop}[Distributividade da União]

  Para quaisquer conjuntos (A), (B) e (C),

  [
    A \cup (B \cap C)
    =
    (A \cup B)
    \cap
    (A \cup C).
  ]

\end{prop}

\begin{prop}[Distributividade da Interseção]

  Para quaisquer conjuntos (A), (B) e (C),

  [
    A \cap (B \cup C)
    =
    (A \cap B)
    \cup
    (A \cap C).
  ]

\end{prop}

\begin{prop}[Leis do Conjunto Vazio]

  Para todo conjunto (A),

  [
    A \cup \varnothing = A,
  ]

  e

  [
    A \cap \varnothing = \varnothing.
  ]

\end{prop}

\begin{prop}[Leis do Complemento]

  Para todo conjunto (A \subseteq U),

  [
    A \cup A^{c}
    =
    U,
  ]

  e

  [
    A \cap A^{c}
    =
    \varnothing.
  ]

\end{prop}

\begin{prop}[Primeira Lei de De Morgan]

  Para quaisquer conjuntos (A) e (B),

  [
    (A \cup B)^c
    =
    A^c \cap B^c.
  ]

\end{prop}

\begin{proof}

  Seja (x\in U).

  Então

  [
    x \in (A\cup B)^c
  ]

  se, e somente se,

  [
    x \notin A\cup B.
  ]

  Pela definição de união,

  [
    \neg(x\in A \lor x\in B).
  ]

  Pela Lei de De Morgan da lógica proposicional,

  [
    (x\notin A)
    \land
    (x\notin B).
  ]

  Logo,

  [
    x \in A^c \cap B^c.
  ]

  Portanto,

  [
    (A\cup B)^c
    =
    A^c\cap B^c.
  ]

\end{proof}

\begin{prop}[Segunda Lei de De Morgan]

  Para quaisquer conjuntos (A) e (B),

  [
    (A \cap B)^c
    =
    A^c \cup B^c.
  ]

\end{prop}

\end{document}
